% GENERATED FILE. Edit canonical lessons, then run npm run generate:books.
% canonical-source-hash: fnv1a64:a7023acfbc95b8d9
% canonical-lessons: TE-C17-madhyaahnam-ardharaatri, TE-S138-letter-gha, TE-S155-letter-ttha, TE-S156-script-recall, TE-S160-letter-pha, TE-S161-script-recall-four

\chapter{Noon and Midnight}
\label{ch:te-noon-and-midnight}
\hlchaptermodality{17}

\begin{chapteropening}
\textbf{By the end of this chapter:} \emph{I can say noon and midnight in Telugu and explain why one is \textquotedblleft{}middle\textquotedblright{} and the other \textquotedblleft{}half.\textquotedblright{}}

Say \te{మధ్యాహ్నం} and \te{అర్ధరాత్రి}, then separate madhya (\textquotedblleft{}middle\textquotedblright{}) from ardha (\textquotedblleft{}half\textquotedblright{}) — two different Sanskrit roots where Kannada uses one.
\end{chapteropening}

\section[madhyāhnam ardharātri]{\te{మధ్యాహ్నం}, \te{అర్ధరాత్రి} — the same noon word, a different midnight word}

\label{lesson:TE-C17-madhyaahnam-ardharaatri}



\begin{quote}
Telugu and Kannada share the exact same Sanskrit word for noon — but then quietly part ways on midnight.
\end{quote}

\begin{cousinweb}
\textbf{\te{మధ్యాహ్నం}} (\emph{madhyāhnam}) = \textquotedblleft{}\textbf{noon, midday}\textquotedblright{} — the same Sanskrit \emph{tatsama} compound as Kannada's \emph{madhyāhna}: \textbf{\te{మధ్య}} (\emph{madhya}, \textquotedblleft{}\textbf{middle}\textquotedblright{}) + \textbf{\te{అహ్న}} (\emph{ahna}, \textquotedblleft{}\textbf{day}\textquotedblright{}), plus Telugu's own \textbf{-\te{ం}} (\emph{anusvāra}) ending. Same root, same meaning, same everyday register in both languages.
\end{cousinweb}

\begin{cousinweb}
\textbf{\te{అర్ధరాత్రి}} (\emph{ardharātri}) = \textquotedblleft{}\textbf{midnight}\textquotedblright{} — but here Telugu reaches for a \textbf{different} Sanskrit root: \textbf{\te{అర్ధ}} (\emph{ardha}, \textquotedblleft{}\textbf{half}\textquotedblright{}), not \emph{madhya} (\textquotedblleft{}middle\textquotedblright{}). Paired with \textbf{\te{రాత్రి}} (\emph{rātri}, \textquotedblleft{}\textbf{night}\textquotedblright{}), it literally means \textquotedblleft{}\textbf{half-night}\textquotedblright{} — structurally the same shape as Hindi's colloquial \emph{ādhī rāt} (\textquotedblleft{}half night\textquotedblright{}) and Latin's transparent \emph{media nox} pattern, but built from Sanskrit's own dedicated \textquotedblleft{}half\textquotedblright{} word rather than \textquotedblleft{}middle.\textquotedblright{}
\end{cousinweb}

\begin{grammarlens}[title={Grammar Lens: Be honest: two different \textquotedblleft{}middle-ish\textquotedblright{} ideas, side by side}]
Don't assume \emph{madhya} and \emph{ardha} are the same word wearing different clothes — they're genuinely \textbf{separate Sanskrit roots} (\textquotedblleft{}middle\textquotedblright{} vs. \textquotedblleft{}half\textquotedblright{}) that happen to land on the same meaning here, the way English \textquotedblleft{}midday\textquotedblright{} and \textquotedblleft{}half the day gone\textquotedblright{} would both point at noon without sharing a root. Telugu's noon word keeps Kannada's \textquotedblleft{}middle\textquotedblright{} logic; its midnight word switches to \textquotedblleft{}half\textquotedblright{} logic instead.
\end{grammarlens}

\subsection*{Your turn}
Say these aloud:

\begin{itemize}
\raggedright
  \item \textquotedblleft{}madhyāhnam\textquotedblright{} — noon, same root as Kannada
  \item \textquotedblleft{}ardharātri\textquotedblright{} — midnight, literally \textquotedblleft{}half-night\textquotedblright{}
  \item the contrast — \textquotedblleft{}madhya\textquotedblright{} (middle) for noon, \textquotedblleft{}ardha\textquotedblright{} (half) for midnight — two different Sanskrit roots
\end{itemize}

\begin{tcolorbox}[breakable,colback=teal!4,colframe=teal!35!black,title={Before you move on}]
Does Telugu's noon word match Kannada's? (\textbf{Yes} — both \emph{madhyāhna(m)}, from \emph{madhya}, \textquotedblleft{}middle.\textquotedblright{}) Does Telugu's midnight word use the same root? (\textbf{No} — \emph{ardharātri} uses \textbf{ardha}, \textquotedblleft{}half,\textquotedblright{} a different Sanskrit root.) What does \emph{ardharātri} literally mean? (\textquotedblleft{}\textbf{Half-night}\textquotedblright{} — the same shape as Hindi's \emph{ādhī rāt}.)
\end{tcolorbox}
\section[gha]{\te{ఘ} — one character, met inside words you already say}

\label{lesson:TE-S138-letter-gha}



\begin{quote}
Before the new one: \te{ఖ} — what does it do? And one from further back: \te{గ}?

One character this time — and you have been saying it for pages without knowing which mark on the page it was.
\end{quote}

\subsection*{Script you'll notice: \te{ఘ}}
\textbf{\te{ఘ}} — \emph{gha}.

It is a \textbf{consonant}, and in this script a consonant is never bare: it comes with an \emph{a} already in it. So it is not \emph{gh}, it is \textbf{gha}.

The same trick as last time, one row down. \textbf{\te{గ}} is \emph{ga}; \textbf{\te{ఘ}} is \emph{ga} with the breath let out behind it.

That is now two breathed pairs — \textbf{\te{క}/\te{ఖ}} and \textbf{\te{గ}/\te{ఘ}} — and they behave identically. When you meet a Telugu letter you do not know, it is worth asking first whether it is a letter you \emph{do} know wearing extra breath.

You already say this one, and the breathed letter is inside it:

\begin{itemize}
\raggedright
  \item \textbf{\te{మాఘం}} \emph{Māghaṁ} — the eleventh of the twelve months
\end{itemize}

\subsection*{Writing: \te{ఘ} — copy what you see}
Put your pen on \te{ఘ} and follow its line. Copy the shape you can see — slowly, and larger than it is printed.

\begin{quote}
This book does not yet tell you \textbf{where to start the character or which way to travel}. That is a real thing, taught with real variation from school to school, and it is not written down here until it can be written down with a source. Copying what is in front of you needs no such source, and it is how the shape gets into your hand in the meantime.
\end{quote}

\subsection*{Your turn}
\begin{itemize}
\raggedright
  \item \emph{Look:} at these two, and say which one has \te{ఘ} and which has \te{గ}
\end{itemize}

\begin{quote}
\te{మాఘం}  ·  \te{బాగా}
\end{quote}

\begin{itemize}
\raggedright
  \item \emph{Trace it:} \te{ఘ} three times, saying \emph{gha} and letting the breath out as you finish each one
  \item \emph{Look:} back at any page of this chapter and find \te{ఘ} once more
\end{itemize}

\begin{tcolorbox}[breakable,colback=teal!4,colframe=teal!35!black,title={Before you move on}]
Which character is this — \te{ఘ}? What sound does it carry, and which plain letter is it built on? (\emph{\textbf{gha}}, built on \te{గ}.)

One more, from much earlier: \te{య} — what does it do?
\end{tcolorbox}
\section[ṭha]{\te{ఠ} — the curled t, with the breath let out}

\label{lesson:TE-S155-letter-ttha}



\begin{quote}
Before the new one: \te{ట} — what sound does it carry, and where does your tongue go for it?

Today's character is that one with a puff of air after it.
\end{quote}

\subsection*{Script you'll notice: \te{ఠ}}
\textbf{\te{ఠ}} — \emph{ṭha}.

It is a \textbf{consonant}, and in this script a consonant is never bare: it comes with an \emph{a} already in it. So it is not \emph{ṭh}, it is \textbf{ṭha}.

\begin{itemize}
\raggedright
  \item \textbf{\te{జ్యేష్ఠం}} \emph{jyēṣṭham} — Jyeshtha, the third month of the year (jyēṣṭham — \textquotedblleft{}the eldest\textquotedblright{})
\end{itemize}

\textbf{\te{ట} and \te{ఠ} are the same tongue position with and without the breath}, the pair this script marks and English does not. The romanization writes them \emph{ṭ} and \emph{ṭha}, which is a difference of two letters for a difference you can hear; on the page they are two separate shapes, and that is the more reliable signal.

This one is rare. It is worth a page because a month name needs it, and a reader who meets an unfamiliar shape in a familiar list has no way to tell whether the word or the printing is at fault.

\subsection*{Writing: \te{ఠ} — copy what you see}
Put your pen on \te{ఠ} and follow its line. Copy the shape you can see — slowly, and larger than it is printed. Then write \te{ట} beside it and look at the two.

\begin{quote}
This book does not yet tell you \textbf{where to start the character or which way to travel}. That is a real thing, taught with real variation from school to school, and it is not written down here until it can be written down with a source. Copying what is in front of you needs no such source, and it is how the shape gets into your hand in the meantime.
\end{quote}

\subsection*{Your turn}
\begin{itemize}
\raggedright
  \item \emph{Look:} at the month name below and find \te{ఠ} near its end
\end{itemize}

\begin{quote}
\te{జ్యేష్ఠం}
\end{quote}

\begin{itemize}
\raggedright
  \item \emph{Trace it:} \te{ఠ} three times, saying \emph{ṭha} as you finish each one
  \item \emph{Say it:} which of \te{ట} and \te{ఠ} has the breath after it
\end{itemize}

\begin{tcolorbox}[breakable,colback=teal!4,colframe=teal!35!black,title={Before you move on}]
Which character is this — \te{ఠ}? What sound does it carry? (\emph{\textbf{ṭha}}.) What is the only difference between it and \te{ట}? (The breath after it.)
\end{tcolorbox}
\section[āru akṣarālu]{\te{ఆరు} \te{అక్షరాలు} — six characters, from cold}

\label{lesson:TE-S156-script-recall}



\begin{quote}
Close the lessons before this one. Say the greeting and the thank-you word, and name the character each of them carries that you met most recently.
\end{quote}

\subsection*{Script: the six, side by side}
\begin{quote}
\te{మ} · \te{ధ} · \te{ట} · \te{ో} · \te{ొ} · \te{ఠ}
\end{quote}

Four are letters and two are signs. \textbf{A letter can stand where a word begins; a sign has to hang on a letter}, and telling which is which is the first question to ask of an unfamiliar shape.

\noindent
\begin{tabularx}{\linewidth}{@{}>{\raggedright\arraybackslash}X>{\raggedright\arraybackslash}X@{}}
\toprule
\textbf{letters} & \textbf{signs} \\
\midrule
\te{మ} \te{ధ} \te{ట} \te{ఠ} & \te{ో} \te{ొ} \\
\bottomrule
\end{tabularx}

\subsection*{Your turn}
\begin{itemize}
\raggedright
  \item \emph{Look:} at these words and name every character in them you have been taught
\end{itemize}

\begin{quote}
\te{నమస్కారం}  ·  \te{ధన్యవాదములు}  ·  \te{ఏమిటి}  ·  \te{తొమ్మిది}
\end{quote}

\begin{itemize}
\raggedright
  \item \emph{Say it:} which of the six is the one you have seen most often
  \item \emph{Say it:} which two of the six cannot start a word
  \item \emph{Trace it:} each of the four letters once, saying its sound
\end{itemize}

\begin{tcolorbox}[breakable,colback=teal!4,colframe=teal!35!black,title={Before you move on}]
Which two of the six are signs rather than letters? (\textbf{\te{ో}} and \textbf{\te{ొ}}.) Which pair is one tongue position with and without the breath? (\textbf{\te{ట}} and \textbf{\te{ఠ}}.) Which of the two \emph{o} signs is the long one? (\textbf{\te{ో}}.)
\end{tcolorbox}
\section[pha]{\te{ఫ} — \te{ప}, with the breath let out}

\label{lesson:TE-S160-letter-pha}



\begin{quote}
A page ago you met \te{ట} and \te{ఠ} — one tongue position, with and without a puff of air after it. Say both.

Now say \te{ప}. Today's character does to \te{ప} what \te{ఠ} did to \te{ట}.
\end{quote}

\subsection*{Script you'll notice: \te{ఫ}}
\textbf{\te{ఫ}} — \emph{pha}.

It is a \textbf{consonant}, and in this script a consonant is never bare: it comes with an \emph{a} already in it. So it is not \emph{ph}, it is \textbf{pha}.

\begin{itemize}
\raggedright
  \item \textbf{\te{ఫాల్గుణం}} \emph{phālguṇam} — Phalguna, the twelfth month of the year
\end{itemize}

Hold a hand in front of your mouth. Say \te{ప}: the air stays put. Say \te{ఫ}: the air comes out. That puff is the whole difference, and Telugu gives it a shape of its own rather than leaving you to guess from the neighbours.

\noindent
\begin{tabularx}{\linewidth}{@{}>{\raggedright\arraybackslash}X>{\raggedright\arraybackslash}X@{}}
\toprule
\textbf{no breath} & \textbf{breath} \\
\midrule
\te{ట} & \te{ఠ} \\
\te{ప} & \te{ఫ} \\
\bottomrule
\end{tabularx}

\textbf{Read the romanization carefully here}: \emph{ph} is one sound with air behind it, not the \emph{f} that the same two English letters make in \emph{phone}. Two letters on the romanized line, one character on the Telugu line, one puff in the mouth.

This character is rare on these pages. It earns a lesson because the last of the twelve months needs it, and an unfamiliar shape inside a list you already know by heart leaves a reader unable to tell whether the word or the printing is at fault.

\subsection*{Writing: \te{ఫ} — copy what you see}
Put your pen on \te{ఫ} and follow its line. Copy the shape you can see — slowly, and larger than it is printed. Then write \te{ప} beside it and look at the two.

\begin{quote}
This book does not yet tell you \textbf{where to start the character or which way to travel}. That is a real thing, taught with real variation from school to school, and it is not written down here until it can be written down with a source. Copying what is in front of you needs no such source, and it is how the shape gets into your hand in the meantime.
\end{quote}

\subsection*{Your turn}
\begin{itemize}
\raggedright
  \item \emph{Look:} at the month name below and find \te{ఫ} at its front
\end{itemize}

\begin{quote}
\te{ఫాల్గుణం}
\end{quote}

\begin{itemize}
\raggedright
  \item \emph{Trace it:} \te{ఫ} three times, saying \emph{pha} as you finish each one
  \item \emph{Say it:} \te{ప} and \te{ఫ} one after the other with a hand in front of your mouth
\end{itemize}

\begin{tcolorbox}[breakable,colback=teal!4,colframe=teal!35!black,title={Before you move on}]
Which character is this — \te{ఫ}? What sound does it carry? (\emph{\textbf{pha}}.) Which month begins with it? (\textbf{\te{ఫాల్గుణం}}.) Name the other pair in this book that differs only by the breath. (\textbf{\te{ట}} and \textbf{\te{ఠ}}.)
\end{tcolorbox}
\section[nālugu akṣarālu]{\te{నాలుగు} \te{అక్షరాలు} — four characters, from cold}

\label{lesson:TE-S161-script-recall-four}



\begin{quote}
Close the lessons before this one. Say how you are, and say the two days at the end of the week. Name the character each of those answers starts on.
\end{quote}

\subsection*{Script: the four, and what each one stands next to}
\begin{quote}
\te{బ} · \te{ష} · \te{శ} · \te{ఫ}
\end{quote}

Not one of these four arrived on its own. Each stands beside a character you had already been reading for chapters, and what it adds is a single change you can hear:

\noindent
\begin{tabularx}{\linewidth}{@{}>{\raggedright\arraybackslash}X>{\raggedright\arraybackslash}X>{\raggedright\arraybackslash}X@{}}
\toprule
\textbf{already yours} & \textbf{new} & \textbf{what changed} \\
\midrule
\te{ప} \emph{pa} & \te{బ} \emph{ba} & the voice switched on \\
\te{ప} \emph{pa} & \te{ఫ} \emph{pha} & a puff of air came after it \\
\te{స} \emph{sa} & \te{ష} \emph{ṣa} & the tongue curled back \\
\te{స} \emph{sa} & \te{శ} \emph{śa} & the tongue went flat against the front of the roof \\
\bottomrule
\end{tabularx}

\textbf{Two of the four hang off \te{ప} and two off \te{స}.} A reader who holds those two anchors in mind has four fewer shapes to memorise cold, because each of the four is one adjustment away from something already known.

\subsection*{Your turn}
\begin{itemize}
\raggedright
  \item \emph{Look:} at these words and name every character in them you have been taught
\end{itemize}

\begin{quote}
\te{బాగా}  ·  \te{సంతోషం}  ·  \te{శనివారం}  ·  \te{ఫాల్గుణం}
\end{quote}

\begin{itemize}
\raggedright
  \item \emph{Say it:} which of the four you have seen most often on these pages
  \item \emph{Say it:} which two of the four are s-sounds
  \item \emph{Trace it:} each of the four once, saying its sound
\end{itemize}

\begin{tcolorbox}[breakable,colback=teal!4,colframe=teal!35!black,title={Before you move on}]
Which of the four differs from \te{ప} by a puff of air? (\textbf{\te{ఫ}}.) Which differs from \te{ప} by the voice alone? (\textbf{\te{బ}}.) Which of \te{ష} and \te{శ} has the tongue curled back? (\textbf{\te{ష}}.) And how many s letters does Telugu print in all, counting the one you had before these two? (\textbf{Three}.)
\end{tcolorbox}
